\section{Expert 的数量与结构设计}

确定了路由方式后，下一个问题是：应该有多少个 expert？每个 expert 应该多大？

\subsection{细粒度 Expert（Fine-Grained Experts）}

\begin{figure}[H]
\centering
\includegraphics[width=0.95\textwidth]{slides-images/slide-018.jpg}
\caption{从标准 MoE 到细粒度 Expert 的演进：将大 expert 切分为更小的 expert，同时增加数量\protect\footnotemark}
\end{figure}
\footnotetext{Slides 第19页。}

早期 MoE 的逻辑演进：

\begin{enumerate}
    \item \textbf{基础 MoE}：将 Dense FFN 复制为多份，每份与原始 FFN 大小相同，Top-2 路由 $\rightarrow$ 激活参数翻倍
    \item \textbf{更多 experts 更好}：实验表明增加 expert 数量能持续降低 loss
    \item \textbf{但不想付出参数代价}：因此将每个 expert 切小——例如将 FFN 的中间维度从 $4d$ 缩小到 $2d$ 或更小，同时 expert 数量翻倍
    \item \textbf{FLOPs 不变}：因为激活的总 FFN 参数量保持不变
\end{enumerate}

\begin{importantbox}{细粒度 Expert 是目前最重要的 MoE 创新之一}
DeepSeek 最早系统性地提出并验证了细粒度 expert 的优势。以 DeepSeek MoE V1 为例：64 个细粒度 routed experts（中间维度为标准的 $1/4$），激活其中 6 个。总 expert 数量增加的同时，每个 expert 更小，带来更灵活的组合能力。这已经被所有主流 MoE 模型采纳。
\end{importantbox}

\subsection{共享 Expert（Shared Experts）}

DeepSeek 的另一个重要创新是\textbf{共享 expert}：在 routed experts 之外，额外设置 1--2 个\textbf{始终激活}的 expert（不经过路由器）。

\begin{knowledgebox}{共享 Expert 的设计动机}
不管输入是什么 token，可能都有一些``通用处理''是必须做的。如果把这些通用计算也交给 routed experts，会浪费路由和参数资源。共享 expert 的作用类似于一个``公共基础设施''，处理所有 token 都需要的通用变换，让 routed experts 专注于差异化处理。
\end{knowledgebox}

\begin{figure}[H]
\centering
\includegraphics[width=0.95\textwidth]{slides-images/slide-020.jpg}
\caption{DeepSeek MoE 消融实验：从 GShard（蓝色）到加入共享 expert（橙色）到细粒度 expert（绿色、红色），性能逐步提升，尤其在 TriviaQA 和 NaturalQuestions 上\protect\footnotemark}
\end{figure}
\footnotetext{Slides 第21页。来自 DeepSeek MoE 论文。}

不过，关于共享 expert 的价值，业界仍有分歧：

\begin{itemize}
    \item DeepSeek V1/V2/V3 使用 2 个共享 expert
    \item OLMoE 的消融实验表明，1 个共享 expert 的增益并不总是显著
    \item 一些中国 LLM 公司曾尝试多个共享 expert，后来又回到 0 或 1 个
    \item 共享 expert 的实际价值可能在于让所有 expert 大小一致，简化系统实现
\end{itemize}

\subsection{OLMoE 消融：更多 experts vs 更大 experts}

\begin{figure}[H]
\centering
\includegraphics[width=0.95\textwidth]{slides-images/slide-021.jpg}
\caption{OLMoE 消融实验：细粒度 experts（8$\rightarrow$32$\rightarrow$64）持续带来增益；共享 expert 的效果则取决于任务\protect\footnotemark}
\end{figure}
\footnotetext{Slides 第22页。来自 OLMoE。}

OLMoE 的系统性消融验证了以下结论：

\begin{itemize}
    \item \textbf{从 8 到 64 个细粒度 experts}：持续降低验证 loss
    \item \textbf{共享 expert}：在某些任务上有帮助，但不如细粒度 expert 那样一致
    \item \textbf{增加 activated experts}：当 expert 更细粒度时，增加 $K$ 是``免费''的（因为每个 expert 更小）
\end{itemize}

\subsection{本章小结}

\begin{itemize}
    \item \textbf{细粒度 expert} 是 MoE 最重要的设计创新——将 expert 切小、数量增多，在不增加 FLOPs 的前提下获得更灵活的组合能力
    \item \textbf{共享 expert} 捕获通用处理需求，但其边际收益仍有争议
    \item \textbf{更多 experts 几乎总是更好}，只要系统能高效支持
\end{itemize}

%% ========== Part 4: 训练挑战 ==========

